\subsection{Local row modules with a semisimple syzygy}

\begin{lemma}[Proved: a semisimple syzygy detects a bounded two-term complex]
\label{CP2:A-local-semisimple-syzygy-two-term}
Let $A$ be a finite-dimensional rational algebra with a unital character
$\chi:A\to\mathbb Q$ whose kernel $J$ is nilpotent, and put
$k={}_\chi\mathbb Q$.
Let $V$ be a finite-dimensional left $A$-module.  Suppose either $V$ is
projective, or there is an exact sequence
\[
 0\longrightarrow k^{\oplus s}\longrightarrow P_{d-1}
 \longrightarrow\cdots\longrightarrow P_0\longrightarrow V
 \longrightarrow0,
 \qquad s\ge1,\quad d\ge1,
\]
in which every $P_i$ is finitely generated projective.  The conclusion
also holds when $V$ itself is semisimple.
If $D$ is a complex of finite-dimensional right $A$-modules in degrees
zero and one, and $D\otimes_A^{\mathbf L}V$ has bounded homology, then
\[
 H_j(D\otimes_A^{\mathbf L}V)=0\qquad(j\ge2).
\]
\end{lemma}
\begin{proof}
The projective case and $V=0$ are immediate.  In the semisimple case,
$V=k^{\oplus s}$ for some $s\ge1$, and boundedness already implies
boundedness of $D\otimes_A^{\mathbf L}k$.
In the displayed exact-sequence case, tensor each short syzygy sequence
with $D$.  For a projective coefficient $P_i$, the tensor is represented
by an ordinary complex in degrees zero and one.  Its homology therefore
vanishes above degree one.  The resulting long exact sequences give
\[
 H_j(D\otimes_A^{\mathbf L}V)
 \cong H_{j-d}(D\otimes_A^{\mathbf L}k)^{\oplus s}
 \qquad(j\ge d+2).
\]
Thus $D\otimes_A^{\mathbf L}k$ has bounded homology in this case too.

For completeness, the minimal-complex argument gives more than a
dimension shift.  Take a bounded-below finite-rank minimal free
resolution $P\to D$.  Its differential matrices lie in $J$.
Consequently $P\otimes_A k$ has zero differentials, and boundedness
implies that $P$ has a highest nonzero term.  If its degree were
$t\ge2$, then $H_t(P)=H_t(D)=0$, so the top differential would be
injective.  Choose a nonzero $x\in A$ with $Jx=0$, using the last
nonzero radical power (or $x=1$ when $J=0$).  Multiplying a free basis
vector of $P_t$ on the right by $x$ gives a nonzero vector killed by
that differential, a contradiction.  Hence $P_j=0$ for every $j\ge2$.
Tensoring this same two-term free resolution with $V$ proves the
assertion, including the degree-two endpoint.
This is the minimal-resolution construction of
Lemma~\ref{CP2:A-local-two-term-derived-amplitude}.
\end{proof}

\begin{corollary}[Proved: arbitrary coefficients for a local square-zero radical]
\label{CP2:A-local-square-zero-two-term}
Under the algebra hypotheses of the preceding lemma, suppose $J^2=0$.
For every finite-dimensional left $A$-module $V$ and every two-term
right complex $D$, boundedness of $D\otimes_A^{\mathbf L}V$ forces
its homology to vanish in degrees at least two.
\end{corollary}
\begin{proof}
Take a minimal projective cover $A^{\oplus t}\twoheadrightarrow V$.
Its kernel $U$ lies in $J^{\oplus t}$, so $JU=0$.
If $U=0$, then $V$ is projective.  Otherwise $U$ is a nonzero
direct sum of copies of the residue module $k$, and the preceding
lemma applies with $d=1$.
\end{proof}

\begin{theorem}[Proved: two-column graph sources with a local semisimple syzygy]
\label{CP2:A-local-semisimple-syzygy-row-graph}
Let $M$ have a proper rectangular minimal band $K$ with two columns,
put $A=\mathbb QM/\mathbb QK$, and let $V$ be its full left
row-augmentation module.  Suppose $A$ has a unital rational character
with nilpotent kernel, and $V$ satisfies one of the coefficient
hypotheses of Lemma~\ref{CP2:A-local-semisimple-syzygy-two-term}.
For a finite nonempty left-$N$, right-$M$ biset, assume the rationally
acyclic opposite corner and natural augmented graph-chain comparison
of Theorem~\ref{CP2:A-two-column-two-outside-graph}.
If the actual collage has bounded rational homology, all its rational
degree-two products vanish.  In particular this holds for $N=R_3$
and every state count.  It also holds for arbitrary $V$ if the
radical of $A$ has square zero.
\end{theorem}
\begin{proof}
The common-image-vertex case in the two-column graph argument gives
an augmentation section on the actual bridge, so every second class
is relative and its products vanish.
Otherwise the two column images have disjoint support.  The injective
projective subcomplex $Ce\otimes e\mathbb QM$ has two-term quotient
$D$ over $A$.  The augmentation sequence and the projective minimal
ideal give the established actual-cofiber comparison
\[
 H_n(Y;\mathbb Q)
   \cong H_{n-2}(D\otimes_A^{\mathbf L}V),\qquad n\ge4.
\]
Here $Y$ is the actual collage with the rationally acyclic opposite
corner collapsed.  These localization steps do not impose any bound
on the number of rows or outside elements.
Boundedness of the actual collage gives boundedness of the tensor.
The lemma, or its square-zero corollary, now gives $H_4(Y;\mathbb Q)=0$.
The actual rational cohomology-ring comparison completes the proof.
\end{proof}

\begin{corollary}[Proved: the original ten-element source $156674$]
\label{CP2:A-ten-source156674-local-graph}
The original positive source with index $156674$ in the complete
order-ten band catalog contributes no $\mathbb CP^2$ collage opposite
$R_3$, for any state count.
\end{corollary}
\begin{proof}
Its minimal band has three rows and two columns, labelled
$e_{00},e_{01},e_{10},e_{11},e_{20},e_{21}$ by $1,\ldots,6$.
The outside elements are $1_M,a,b,c$, labelled $0,7,8,9$.
Every product of two elements of $\{a,b,c\}$ lies in the band.
Thus the contracted algebra is
\[
 A=\mathbb Q1\oplus J,\qquad
 J=\mathbb Qa\oplus\mathbb Qb\oplus\mathbb Qc,
 \qquad J^2=0.
\]
In the row-difference basis $v_1,v_2$, the three actors are
$a,c\mapsto E_{12}$ and $b\mapsto-E_{12}$.
One may either use the square-zero corollary directly, or the explicit
cyclic cover $A\to V$, $x\mapsto xv_2$, whose kernel is
\[
 \mathbb Q(a+b)\oplus\mathbb Q(c-a)\cong k^{\oplus2}.
\]
The exact literal multiplication, row-module and kernel identities
are bound in
\path{converse/scope18/ten_source156674_local_binding.json}.
The graph theorem applies directly to the hypothetical original
collage, whose homology is bounded.
\end{proof}

\begin{remark*}
The last corollary is applied to the original source table.  A map
preserving only second cohomology does not supply bounded higher
homology for its target; no such transfer is used here.
\end{remark*}
